\chapter{ویزای تحصیلی}
\label{ch:visa}
گرفتن ویزا مهم‌ترین مرحله مهاجرت تحصیلی شما و بسته به کشوری که قصد دارید ویزای آن را بگیرید 
ممکن است سخت یا آسان باشد. هر کشوری از قوانین خاص خود برای ویزا پیروی می‌کند و \textbf{به هیچ وجه} نمی‌توان یک 
قاعده مشخص برای ویزای کشور‌های مختلف مشخص کرد. با این حال نکاتی وجود دارد که بین کشور‌هایی که من (نویسنده این متن) با آن 
برخورد داشته‌ام مشترک بوده که در ادامه آن‌ها را ذکر می‌کنم ولی به طور کلی قوانین اخذ ویزای تحصیلی را 
باید از سفارت و گرو‌ه‌ها و منابع مختص آن پیگیری کنید.

\section{مدارک رایج}
\subsection{پاسپورت معتبر}
پاسپورت دارای اعتبار بیش از ۶ ماه (برخی کشور‌ها یک سال) در زمان ثبت درخواست ویزا یکی از رایج‌ترین مدارک مورد نیاز سفارت‌هاست. پاسپورت با 
اعتبار اندک ممکن است منجر به نقص مدرک بشود. درباره دانشجویان پسر که گرفتن پاسپورت برای آن‌ها نیاز به تایید نظام وظیفه دارد 
توصیه می‌شود در زمان تحصیل پاسپورت خود را گرفته و از اعتبار آن مطمئن شوید. 

\subsection{نامه پذیرش}
جهت اثبات قصد تحصیل در کشور مقصد نامه‌پذیرش نیز یکی از مهم‌ترین مدارک مورد نیاز برای ارائه به سفارت است. این نامه که توسط
دانشگاه/موسسه مقصد صادر می‌شود نشان می‌دهد شما توسط آن‌ها تحت شرایط مشخص پذیرفته شده‌اید و برای یک دوره تحصیل با طول تقریبی مشخص
و در شرایط مشخص در آن دانشگاه تحصیل خواهید کرد. 

\subsection{تمکن مالی}
جهت اثبات اینکه شما توانایی تامین هزینه‌های حین تحصیل خود را دارید، سفارت از شما مدارکی می‌خواهد که نشان دهنده توانایی مالی شما باشد. 
تمکن مالی باید به صورت \lr{liquid asset} باشد، با توجه به تجربیات من جدول زیر نشان می‌دهد که کدام موارد می‌تواند به عنوان
تمکن ارائه شود. 
\begin{table}[h]
\centering
\caption{انواع دارایی قابل ارائه به عنوان تمکن مالی}
\label{tab:financial-capability}
\begin{tabular}{|p{4cm}|p{3cm}|p{7cm}|}
\hline
\textbf{نوع دارایی} & \textbf{قابلیت ارائه} & \textbf{توضیحات} \\
\hline
موجودی حساب بانکی & قابل قبول & موجودی انواع حساب بانکی خودتان یا اسپانسر که قابل برداشت باشد مورد قبول است. ولی گاهی باید مدارکی ارائه شود که منبع پول یا قرضی نبودن موجودی را اثبات کند \\
\hline
بورس، سهام و صندوق‌های سرمایه‌گذاری & با احتیاط & در \textbf{برخی موارد ولی نه همیشه} قابل قبول است، اما به دلیل نوسان قیمت و زمان لازم برای نقد شدن، ممکن است سفارت آن را به اندازه موجودی نقدی معتبر نداند. \\
\hline
طلا، سکه، زمین، ملک، خودرو، رمز ارز و غیره & غیرقابل قبول & اسناد این دارایی‌ها به عنوان نشان دهنده پایداری مالی خانواده یا وابستگی به کشور قابل ارائه است ولی به عنوان تمکن قابل پذیرش نیست زیرا دارایی قابل انتقال محسوب نمی‌شود \\
\hline
فیش حقوقی یا درآمد ماهانه & مکمل & فیش حقوقی به تنهایی تمکن مالی محسوب نمی‌شود، اما می‌تواند به عنوان مدرکی برای اثبات منبع درآمد و پشتیبانی مالی ارائه شود. \\
\hline
\end{tabular}
\end{table}

با این حال قواعد ارائه مدارک تمکن مالی برای کشور‌های مختلف متفاوت است، بعضی از کشور‌ها از شما می‌خواهند هزینه سال اول را نشان داده 
و بعضی هزینه کل دوره تحصیل را خواهند خواست، همچنین میزان ماندگاری پول، گردش حساب و موارد مشابه از جمله مواردی است که درباره 
کشور‌های مختلف متفاوت است.

\begin{tipbox}
    برگه تمکن مالی توسط بانک‌های مختلف تحت قوانین و فرمت‌های نسبتا مشابهی صادر می‌شود، مثلا بعضی بانک‌ها (مثلا سامان) به شما
    اجازه می‌دهند تمکن را در برای هر بازه زمانی دلخواه بگیرید ولی برخی دیگر فقط اجازه انتخاب بازه‌های مشخصی را به شما می‌دهند. 
    این فرمت‌ها گاهی در طول زمان تغییر می‌کند ولی با این حال نمونه برگه تمکن بانک‌ها را در فایل 
    \href{https://github.com/farbodbj/ApplyWay/blob/master/assets/tamakon-example.pdf}{نمونه تمکن‌بانکهای مختلف \footnote{این فایل بیش از سه سال پیش تهیه شده به همه موارد آن اطمینان نکنید} - گروه کانادا دریم}
    ببینید.
\end{tipbox}


\section{اطلاعات مربوط به هر کشور}
در این بخش اطلاعات مربوط به ویزای کشور‌هایی که مقصد متداول‌تری برای دانشجویان محسوب می‌شود آورده می‌شود. تلاش شده به جای ارائه مطالب به صورت 
یکجا لینک به منابع معتبر آورده شود که در زمان‌های مختلف معتبر باشد. 


\subsection{کانادا}
فرایند ویزای کانادا به طور خاص برای ایرانیان شامل جزئیات و نکات زیادی می‌شود. از نوشتن متن \lr{Client Info} برای توجیه دلیل 
تحصیل در یک مقطع خاص تا تمکن مالی با شرایط مناسب برای تامین هزینه‌های تحصیل. ذکر کردن جزئیات در این باره از حوصله این راهنما خارج است 
و برای کانادا شما را به دفترچه راهنمای به خصوص آن ارجاع می‌دهیم که توسط گروه اپلای کاما تهیه شده ارجاع می‌دهیم. این راهنما بسیار کامل است و 
نکات زیادی را در بر میگیرد با این حال عضویت در گروه‌های مرتبط، سوال کردن درباره جدیدترین آپدیت‌ها و قوانین کاملا ضروری است. 

در ادامه لینک فایل‌های راهنمای مذکور را با شما به اشتراک می‌گذاریم. 

\href{https://github.com/farbodbj/ApplyWay/blob/master/assets/canada-visa/study-permit-temporary-stay.pdf}{راهنمای ویزا تحصیلی کانادا - گروه کاما}

\href{https://github.com/farbodbj/ApplyWay/blob/master/assets/canada-visa/tracking-your-application.pdf}{بررسی وضعیت پرونده پس از ثبت - گروه کانادا دریم}

\subsection{ایتالیا}
کامل شود

\subsection{اسپانیا}
کامل شود


\section{لینک‌های مفید}
\label{sec:visa-useful-links}
در جدول زیر پیوندهای رسمی مربوط به ویزای تحصیلی یا مهاجرت چند کشور مقصد رایج برای دانشجویان ایرانی
آمده است. تلاش شده علاوه بر آمریکای شمالی، اروپا و اقیانوسیه، مقصدهای شرق آسیا، جنوب‌شرق آسیا،
خاورمیانه و قفقاز هم پوشش داده شوند. این پیوندها نقطه شروع پیگیری قوانین روز هستند؛ آدرس دقیق
سفارت یا مرکز انگشت‌نگاری ممکن است بسته به محل اقامت شما (ایران یا کشور ثالث) فرق کند.

\begin{warningbox}
  بعضی کشور‌ها برای ایرانیان سفارت فعال در تهران ندارند و پرونده از طریق کشور ثالث یا مراکز
  خدماتی مثل \lr{VFS} / \lr{TLS} بررسی می‌شود. همیشه منبع رسمی همان کشور مقصد را ملاک قرار دهید.
\end{warningbox}

\begin{warningbox}
  بعضی کشورها برای ایرانیان سفارت فعال در تهران ندارند و پرونده از طریق کشور ثالث یا سفارتخانه‌های آکرودیته بررسی می‌شود. همواره آخرین اطلاعات را از وب‌سایت رسمی دولت یا سفارت کشور مقصد بررسی کنید.
\end{warningbox}

\begin{longtable}{|p{2.4cm}|p{2.2cm}|p{8.4cm}|}
\caption{پیوندهای رسمی ویزا و سفارتخانه‌های مرتبط با ایران}
\label{tab:visa-official-links}\\
\hline
\textbf{کشور} & \textbf{منطقه} & \textbf{پیوند رسمی و سفارت} \\
\hline
\endfirsthead
\hline
\textbf{کشور} & \textbf{منطقه} & \textbf{پیوند رسمی و سفارت} \\
\hline
\endhead
\hline
\endfoot

کانادا & آمریکای شمالی &
\href{https://www.canada.ca/en/services/immigration-citizenship.html}{\lr{IRCC}}\\
& & {\small سفارت کانادا در تهران فعال نیست؛ درخواست‌ها معمولاً از طریق آنکارا یا ابوظبی بررسی می‌شوند.}\\
\hline

آمریکا & آمریکای شمالی &
\href{https://travel.state.gov/content/travel/en/us-visas.html}{\lr{U.S. Department of State -- Visas}}\\
& & {\small سفارت آمریکا در تهران وجود ندارد؛ ایرانیان معمولاً از طریق آنکارا، ایروان، دبی یا سایر سفارتخانه‌ها اقدام می‌کنند.}\\
\hline

بریتانیا & اروپا &
\href{https://www.gov.uk/student-visa}{\lr{GOV.UK -- Student Visa}}\\
& & {\small \href{https://www.gov.uk/world/organisations/british-embassy-tehran}{\lr{British Embassy Tehran}}}\\
\hline

آلمان & اروپا &
\href{https://www.auswaertiges-amt.de/en/visa-service}{\lr{German Federal Foreign Office -- Visa}}\\
& & {\small \href{https://teheran.diplo.de/}{\lr{Embassy of Germany in Tehran}}}\\
\hline

فرانسه & اروپا &
\href{https://france-visas.gouv.fr/}{\lr{France-Visas}}\\
& & {\small \href{https://ir.ambafrance.org/}{\lr{Embassy of France in Tehran}}}\\
\hline

ایتالیا & اروپا &
\href{https://vistoperitalia.esteri.it/}{\lr{Visa for Italy}}\\
& & {\small \href{https://ambteheran.esteri.it/}{\lr{Embassy of Italy in Tehran}}}\\
\hline

اسپانیا & اروپا &
\href{https://www.exteriores.gob.es/Consulados/teheran/en/ServiciosConsulares/Paginas/Consular/Visados.aspx}{\lr{Spain Visa Information}}\\
& & {\small \href{https://www.exteriores.gob.es/Consulados/teheran}{\lr{Embassy of Spain in Tehran}}}\\
\hline

هلند & اروپا &
\href{https://ind.nl/en}{\lr{Immigration and Naturalisation Service (IND)}}\\
& & {\small \href{https://www.netherlandsworldwide.nl/countries/iran/about-us/embassy-in-tehran}{\lr{Embassy of the Netherlands in Tehran}}}\\
\hline

سوئد & اروپا &
\href{https://www.migrationsverket.se/English.html}{\lr{Swedish Migration Agency}}\\
& & {\small \href{https://www.swedenabroad.se/en/embassies/iran-tehran/}{\lr{Embassy of Sweden in Tehran}}}\\
\hline

اتریش & اروپا &
\href{https://www.migration.gv.at/en/}{\lr{Migration Austria}}\\
& & {\small \href{https://www.bmeia.gv.at/oeb-teheran}{\lr{Embassy of Austria in Tehran}}}\\
\hline

فنلاند & اروپا &
\href{https://migri.fi/en/home}{\lr{Finnish Immigration Service (Migri)}}\\
& & {\small \href{https://finlandabroad.fi/web/irn}{\lr{Embassy of Finland in Tehran}}}\\
\hline

ایرلند & اروپا &
\href{https://www.irishimmigration.ie/}{\lr{Irish Immigration Service}}\\
& & {\small \href{https://www.ireland.ie/en/turkey/ankara/}{\lr{Embassy of Ireland in Ankara (accredited to Iran)}}}\\
\hline

پرتغال & اروپا &
\href{https://vistos.mne.gov.pt/en/}{\lr{Portugal Visa Portal}}\\
& & {\small \href{https://teerao.embaixadaportugal.mne.gov.pt/}{\lr{Embassy of Portugal in Tehran}}}\\
\hline

استرالیا & اقیانوسیه &
\href{https://immi.homeaffairs.gov.au/}{\lr{Department of Home Affairs}}\\
& & {\small \href{https://iran.embassy.gov.au/}{\lr{Australian Embassy Tehran}}}\\
\hline

نیوزیلند & اقیانوسیه &
\href{https://www.immigration.govt.nz/}{\lr{Immigration New Zealand}}\\
& & {\small نمایندگی دیپلماتیک ایران تحت پوشش سفارت نیوزیلند در آنکارا است.}\\
\hline

ژاپن & شرق آسیا &
\href{https://www.mofa.go.jp/j_info/visit/visa/index.html}{\lr{MOFA Japan Visa}}\\
& & {\small \href{https://www.ir.emb-japan.go.jp/}{\lr{Embassy of Japan in Tehran}}}\\
\hline

کره جنوبی & شرق آسیا &
\href{https://www.visa.go.kr/}{\lr{Korea Visa Portal}}\\
& & {\small \href{https://overseas.mofa.go.kr/ir-en/index.do}{\lr{Embassy of the Republic of Korea in Tehran}}}\\
\hline

چین & شرق آسیا &
\href{https://www.visaforchina.cn/}{\lr{Chinese Visa Application Service}}\\
& & {\small \href{https://ir.china-embassy.gov.cn/}{\lr{Embassy of China in Tehran}}}\\
\hline

تایوان & شرق آسیا &
\href{https://www.boca.gov.tw/}{\lr{Bureau of Consular Affairs}}\\
& & {\small \href{https://www.roc-taiwan.org/ir_en/}{\lr{Taipei Economic and Cultural Office}}}\\
\hline

هنگ‌کنگ & شرق آسیا &
\href{https://www.immd.gov.hk/eng/services/visas/index.html}{\lr{Hong Kong Immigration Department}}\\
& & {\small خدمات از طریق اداره مهاجرت هنگ‌کنگ ارائه می‌شود.}\\
\hline

سنگاپور & جنوب‌شرق آسیا &
\href{https://www.ica.gov.sg/enter-transit-depart/entering-singapore/visa_requirements}{\lr{Singapore Visa Requirements}}\\
& & {\small سنگاپور سفارت مقیم در تهران ندارد.}\\
\hline

مالزی & جنوب‌شرق آسیا &
\href{https://www.imi.gov.my/}{\lr{Immigration Department of Malaysia}}\\
& & {\small \href{https://www.kln.gov.my/web/irn_tehran}{\lr{Embassy of Malaysia in Tehran}}}\\
\hline

هند & جنوب آسیا &
\href{https://indianvisaonline.gov.in/}{\lr{Indian Visa Online}}\\
& & {\small \href{https://www.indianembassytehran.gov.in/}{\lr{Embassy of India in Tehran}}}\\
\hline

ترکیه & غرب آسیا &
\href{https://www.evisa.gov.tr/en/}{\lr{Turkey e-Visa}}\\
& & {\small \href{https://tehran-emb.mfa.gov.tr/}{\lr{Embassy of Türkiye in Tehran}}}\\
\hline

امارات متحده عربی & خاورمیانه &
\href{https://icp.gov.ae/}{\lr{ICP UAE}}\\
& & {\small \href{https://www.mofa.gov.ae/}{\lr{Embassy of the UAE in Tehran}}}\\
\hline

قطر & خاورمیانه &
\href{https://hukoomi.gov.qa/en/service/apply-for-visa}{\lr{Qatar Visa Services}}\\
& & {\small \href{https://tehran.embassy.qa/}{\lr{Embassy of Qatar in Tehran}}}\\
\hline

روسیه & اوراسیا &
\href{https://electronic-visa.kdmid.ru/}{\lr{Russian e-Visa}}\\
& & {\small \href{https://iran.mid.ru/}{\lr{Embassy of Russia in Tehran}}}\\
\hline

\end{longtable}
\begin{warningbox}
  لیست بالا ممکن است دارای اشکالاتی باشد در صورت مشاهده با ثبت \lr{issue} به ما در بهبود آن کمک کنید. 
\end{warningbox}

